\subsection*{三个阶段}

前序阶段发起第 0 块的装入，用 \co@lf0@ 和 \co@rf0@ 表示。此时还没有上一块可供计算，这两次搬运为后面的循环建立初始状态。

稳定阶段先用 \co@lf1@ 和 \co@rf1@ 发起下一块的装入，再等待当前块就绪，并从 \co@lf0@ 和 \co@rf0@ 读取计算。两对期值对应不同缓冲区，因此下一块的搬运可以与当前计算重叠，并不要求下一块先装入完成。当前计算结束后，\co@swap@ 将下一对变为下一次迭代的当前对。

循环停止发起搬运后，收尾阶段处理最后装入的分块。读取之前，仍需确认它已装入完成。上面的清单用于展示期值轮换，省略了最后一次等待和输出初始化；实际执行累加时必须补齐这两项，也必须像稳定阶段一样保证读取时数据已就绪。
